% ============================================================
% Section 4: SNR Ratio Theory for Graph Neural Networks
% Rigorous information-theoretic framework based on first principles
% Key contribution: Aggregation SNR Ratio theorem and extensions
% ============================================================

\section{Signal-to-Noise Ratio Framework for Structural Enhancement}
\label{sec:snr_theory}

We develop a rigorous information-theoretic framework for understanding when graph structure improves node classification. Our analysis is grounded in first-principles derivation of how message passing affects the signal-to-noise ratio (SNR) of node representations.

\subsection{Problem Setup}
\label{subsec:setup}

\begin{definition}[Contextual Stochastic Block Model with Gaussian Features]
\label{def:csbm}
Consider a graph $G = (V, E)$ with $|V| = n$ nodes. Each node $v \in V$ has:
\begin{itemize}
    \item A binary label $Y_v \in \{+1, -1\}$ with uniform prior
    \item Features $X_v = \mu Y_v + \varepsilon_v$ where $\varepsilon_v \sim \mathcal{N}(0, \sigma^2 I_{d_x})$
\end{itemize}
Edges are generated with \textbf{homophily parameter} $h \in [0, 1]$:
\begin{equation}
    h := \mathbb{P}(Y_u = Y_v \mid (u, v) \in E)
\end{equation}
Define the \textbf{correlation coefficient} $\rho := 2h - 1 \in [-1, 1]$, satisfying:
\begin{equation}
    \mathbb{E}[Y_u \mid (u,v) \in E, Y_v = y] = \rho \cdot y
\end{equation}
\end{definition}

\begin{definition}[Signal-to-Noise Ratio]
\label{def:snr}
For a representation $Z$ with class-conditional distribution $Z \mid Y = y$, the \textbf{SNR} is:
\begin{equation}
    \mathrm{SNR}(Z) := \frac{\|\mathbb{E}[Z \mid Y=+1] - \mathbb{E}[Z \mid Y=-1]\|^2}{\mathrm{tr}(\mathrm{Cov}(Z \mid Y))}
\end{equation}
For the original features under Definition~\ref{def:csbm}:
\begin{equation}
    \mathrm{SNR}_X = \frac{4\|\mu\|^2}{d_x \sigma^2}
\end{equation}
\end{definition}

\subsection{The Aggregation SNR Ratio Theorem}
\label{subsec:snr_ratio}

Our central result characterizes how one-hop aggregation affects the SNR.

\begin{definition}[Mean Aggregation]
\label{def:aggregation}
For node $v$ with degree $k_v := |\mathcal{N}(v)|$, the one-hop mean aggregation is:
\begin{equation}
    A_v := \frac{1}{k_v} \sum_{u \in \mathcal{N}(v)} X_u
\end{equation}
\end{definition}

\begin{theorem}[Aggregation SNR Ratio]
\label{thm:snr_ratio}
Under the CSBM with Gaussian features (Definition~\ref{def:csbm}), assume:
\begin{enumerate}
    \item[(A1)] \textbf{Local independence}: Given $Y_v$, neighbor labels $\{Y_u : u \in \mathcal{N}(v)\}$ are approximately independent (valid in sparse limit $n \to \infty$)
    \item[(A2)] \textbf{Homogeneous degree}: All nodes have degree $k$
\end{enumerate}
Then the SNR after one-hop mean aggregation satisfies:
\begin{equation}
    \kappa := \frac{\mathrm{SNR}_A}{\mathrm{SNR}_X} = \rho^2 \cdot k = (2h-1)^2 \cdot k
\end{equation}
\end{theorem}

\begin{proof}
\textbf{Step 1 (Conditional mean of aggregated features):}
Given $Y_v = y$ and edge $(u,v) \in E$:
\begin{equation}
    \mathbb{E}[X_u \mid (u,v) \in E, Y_v = y] = \mu \cdot \mathbb{E}[Y_u \mid (u,v) \in E, Y_v = y] = \rho \mu y
\end{equation}
Averaging over $k$ neighbors:
\begin{equation}
    \mathbb{E}[A_v \mid Y_v = y] = \rho \mu y
\end{equation}
Thus the class separation after aggregation is:
\begin{equation}
    \|\mathbb{E}[A_v \mid Y_v = +1] - \mathbb{E}[A_v \mid Y_v = -1]\| = 2|\rho| \|\mu\|
\end{equation}

\textbf{Step 2 (Conditional variance of aggregated features):}
Under assumption (A1), neighbor features are conditionally independent given $Y_v$:
\begin{equation}
    \mathrm{Cov}(A_v \mid Y_v) = \frac{1}{k^2} \sum_{u \in \mathcal{N}(v)} \mathrm{Cov}(X_u \mid Y_u) = \frac{1}{k} \sigma^2 I_{d_x}
\end{equation}
Thus:
\begin{equation}
    \mathrm{tr}(\mathrm{Cov}(A_v \mid Y_v)) = \frac{d_x \sigma^2}{k}
\end{equation}

\textbf{Step 3 (SNR ratio):}
\begin{equation}
    \mathrm{SNR}_A = \frac{(2|\rho| \|\mu\|)^2}{d_x \sigma^2 / k} = \rho^2 k \cdot \frac{4\|\mu\|^2}{d_x \sigma^2} = \rho^2 k \cdot \mathrm{SNR}_X
\end{equation}
\end{proof}

\begin{corollary}[Critical Threshold]
\label{cor:threshold}
The aggregation is beneficial ($\kappa > 1$) if and only if:
\begin{equation}
    |\rho| > \frac{1}{\sqrt{k}} \quad \Leftrightarrow \quad h > \frac{1 + 1/\sqrt{k}}{2} \text{ or } h < \frac{1 - 1/\sqrt{k}}{2}
\end{equation}
For $k = 10$: the ``dead zone'' is $h \in (0.34, 0.66)$, matching empirical thresholds $(0.3, 0.7)$.
\end{corollary}

\subsection{From SNR to Classification Error}
\label{subsec:error}

\begin{theorem}[Bayes Error under Aggregation]
\label{thm:bayes_error}
Under the Gaussian class-conditional model, the Bayes-optimal classification error using representation $Z$ is:
\begin{equation}
    P_e(Z) = \Phi\left(-\frac{1}{2}\sqrt{d_x \cdot \mathrm{SNR}(Z)}\right)
\end{equation}
where $\Phi(\cdot)$ is the standard normal CDF.

Consequently, for aggregated features:
\begin{equation}
    P_e(A) = \Phi\left(-\frac{1}{2}\sqrt{d_x \cdot \rho^2 k \cdot \mathrm{SNR}_X}\right)
\end{equation}
\begin{itemize}
    \item $\kappa > 1$: Aggregation \textbf{decreases} Bayes error
    \item $\kappa < 1$: Aggregation \textbf{increases} Bayes error
    \item $\rho = 0$: $P_e(A) = 0.5$ (random guessing)
\end{itemize}
\end{theorem}

\subsection{Multi-Layer SNR Dynamics}
\label{subsec:multilayer}

\begin{theorem}[Multi-Layer SNR Recursion]
\label{thm:multilayer}
Define the $t$-th layer representation $A^{(t)}$ with $A^{(0)} = X$. Under replace aggregation:

\textbf{(a) Without layer noise ($\sigma_\varepsilon^2 = 0$):}
\begin{equation}
    \mathrm{SNR}^{(t)} = (\rho^2 k)^t \cdot \mathrm{SNR}^{(0)} = \kappa^t \cdot \mathrm{SNR}_X
\end{equation}
\begin{itemize}
    \item $\kappa > 1$: SNR grows exponentially (beneficial)
    \item $\kappa < 1$: SNR decays exponentially (harmful, oversmoothing)
    \item $\kappa = 1$: SNR remains constant (critical point)
\end{itemize}

\textbf{(b) With layer noise ($\sigma_\varepsilon^2 > 0$):}
Let $\Delta_t = \|\mathbb{E}[A^{(t)} \mid Y=+1] - \mathbb{E}[A^{(t)} \mid Y=-1]\|$ and $V_t = \mathrm{tr}(\mathrm{Cov}(A^{(t)} \mid Y))$.
\begin{align}
    \Delta_{t+1} &= \rho k \cdot \Delta_t \\
    V_{t+1} &= k \cdot V_t + d_x \sigma_\varepsilon^2
\end{align}
Solving the recursion:
\begin{equation}
    \mathrm{SNR}^{(t)} = \frac{(\rho k)^{2t} \Delta_0^2}{k^t V_0 + d_x \sigma_\varepsilon^2 \frac{k^t - 1}{k - 1}}
\end{equation}
The \textbf{critical layer depth} is:
\begin{equation}
    t^\star = \arg\max_{t \in \mathbb{N}} \mathrm{SNR}^{(t)}
\end{equation}
\end{theorem}

\begin{remark}[Oversmoothing Interpretation]
When $\kappa < 1$ or with layer noise, the SNR eventually decays, causing the ``oversmoothing'' phenomenon where deep GNNs converge to uninformative representations.
\end{remark}

\subsection{GCN Normalization and Effective Degree}
\label{subsec:gcn}

\begin{theorem}[Effective Degree under Normalization]
\label{thm:effective_degree}
For GCN-style aggregation with weights $w_{ij} = 1/\sqrt{d_i d_j}$, define the \textbf{effective degree}:
\begin{equation}
    k_{\mathrm{eff}}(i) := \frac{\left(\sum_{j \in \mathcal{N}(i)} w_{ij}\right)^2}{\sum_{j \in \mathcal{N}(i)} w_{ij}^2}
\end{equation}
Then the local SNR ratio is:
\begin{equation}
    \kappa_i^{\mathrm{(GCN)}} \approx \rho^2 \cdot k_{\mathrm{eff}}(i)
\end{equation}

Special cases:
\begin{itemize}
    \item \textbf{$k$-regular graph}: $k_{\mathrm{eff}} = k$, so $\kappa^{\mathrm{(GCN)}} = \rho^2 k$ (same as unnormalized)
    \item \textbf{Heterogeneous degrees}: $k_{\mathrm{eff}}(i) \ll d_i$ for high-degree nodes, reducing aggregation benefit
\end{itemize}
\end{theorem}

\subsection{Concatenation Preserves Information}
\label{subsec:concat}

\begin{theorem}[SNR under Concatenation]
\label{thm:concat}
For concatenated representation $Z = (X, A)$, assuming $X$ and $A$ are approximately uncorrelated in noise:
\begin{equation}
    \mathrm{SNR}_{\mathrm{concat}} \approx \mathrm{SNR}_X + \mathrm{SNR}_A = (1 + \kappa) \cdot \mathrm{SNR}_X = (1 + \rho^2 k) \cdot \mathrm{SNR}_X
\end{equation}
\end{theorem}

\begin{corollary}[GraphSAGE ADR $\approx 0$ Explained]
\label{cor:sage}
GraphSAGE uses concatenation. When $\rho \approx 0$ (weak homophily):
\begin{itemize}
    \item $\mathrm{SNR}_A \approx 0$
    \item $\mathrm{SNR}_{\mathrm{concat}} \approx \mathrm{SNR}_X$
    \item The model learns to ignore the aggregation branch
    \item ADR $\approx 0$ (no aggregation damage)
\end{itemize}
This explains the observed 200$\times$ ADR difference: GCN (replace, $\kappa \approx 0$) vs GraphSAGE (concat, protected).
\end{corollary}

\subsection{Connection to Kesten-Stigum Threshold}
\label{subsec:ks}

\begin{remark}[Comparison with KS Threshold]
The classical Kesten-Stigum threshold for community detection is $(k-1)\rho^2 = 1$, i.e., $|\rho| = 1/\sqrt{k-1}$.

Our threshold is $k\rho^2 = 1$, i.e., $|\rho| = 1/\sqrt{k}$.

The difference ($k$ vs $k-1$) arises from different mechanisms:
\begin{itemize}
    \item \textbf{KS threshold}: Concerns deep propagation on trees; the effective branching factor is $k-1$ (non-backtracking)
    \item \textbf{Our threshold}: Concerns one-hop SNR improvement; uses all $k$ neighbors for variance reduction
\end{itemize}
Both thresholds are asymptotically equivalent for large $k$, but their mechanisms differ fundamentally.
\end{remark}

\subsection{Summary of Theoretical Contributions}
\label{subsec:summary}

\begin{table}[t]
\centering
\caption{SNR-Based Theoretical Framework: Key Results}
\label{tab:snr_summary}
\begin{tabular}{@{}lcc@{}}
\toprule
\textbf{Result} & \textbf{Formula} & \textbf{Source} \\
\midrule
SNR ratio (1-hop) & $\kappa = \rho^2 k$ & Theorem~\ref{thm:snr_ratio} \\
Bayes error & $P_e = \Phi(-\frac{1}{2}\sqrt{d_x \cdot \mathrm{SNR}})$ & Theorem~\ref{thm:bayes_error} \\
Multi-layer SNR & $\mathrm{SNR}^{(t)} = \kappa^t \cdot \mathrm{SNR}_X$ & Theorem~\ref{thm:multilayer} \\
GCN effective degree & $\kappa_i^{\mathrm{GCN}} = \rho^2 k_{\mathrm{eff}}(i)$ & Theorem~\ref{thm:effective_degree} \\
Concatenation SNR & $\mathrm{SNR}_{\mathrm{concat}} = (1+\kappa) \mathrm{SNR}_X$ & Theorem~\ref{thm:concat} \\
Critical threshold & $|\rho| = 1/\sqrt{k}$ & Corollary~\ref{cor:threshold} \\
\bottomrule
\end{tabular}
\end{table}

\begin{remark}[Scope and Limitations]
\label{rem:scope}
The theoretical results hold under:
\begin{enumerate}
    \item Binary classification with uniform prior
    \item Gaussian features with class-conditional mean
    \item CSBM edge generation
    \item Local independence approximation (valid in sparse limit)
    \item Homogeneous degree assumption (relaxed via $k_{\mathrm{eff}}$)
\end{enumerate}
Extensions to multi-class, non-Gaussian features, and heterogeneous graphs require additional analysis.
\end{remark}

